% GENERATED FILE. Edit canonical lessons, then run npm run generate:books.
% canonical-source-hash: fnv1a64:4b8570a2b3d5d807
% canonical-lessons: TE-C197-vadrangi, TE-C197-racayita, TE-C197-samasya, TE-C197-pariskaram, TE-C197-abhiprayam

\chapter{Carpenter, Writer, Problem, Solution, Opinion}
\label{ch:te-carpenter-writer-problem-solution-opinion}
\hlchaptermodality{197}

\begin{chapteropening}
\textbf{By the end of this chapter:} \emph{I can say a carpenter, a writer, a problem, a solution and an opinion.}

Complete the last lesson of chapter 197: I can say a carpenter, a writer, a problem, a solution and an opinion.
\end{chapteropening}

\section[\texorpdfstring{\te{వడ్రంగి} (vaḍraṅgi)}{వడ్రంగి (vaḍraṅgi)}]{\te{వడ్రంగి} (vaḍraṅgi) — a carpenter}

\label{lesson:TE-C197-vadrangi}



\begin{quote}
Before the new one: say the Telugu for a responsibility, then the Telugu for an application.
\end{quote}

\subsection*{You'll want to know: \te{వడ్రంగి}}
\textbf{\te{వడ్రంగి}} — \emph{vaḍraṅgi} — \textquotedblleft{}a carpenter\textquotedblright{}.

\begin{grammarlens}[title={What you've built}]
One more word for this chapter.
\end{grammarlens}

\subsection*{Your turn}
\begin{itemize}
\raggedright
  \item \emph{Say it:} \emph{vaḍraṅgi}
  \item \emph{Say it:} \emph{vaḍraṅgi}, once more
  \item \emph{Say it:} say \emph{vaḍraṅgi}
  \item \emph{Recall:} say the Telugu for a responsibility, then the Telugu for an application, then say \emph{vaḍraṅgi} again
\end{itemize}

\begin{tcolorbox}[breakable,colback=teal!4,colframe=teal!35!black,title={Before you move on}]
What does \te{వడ్రంగి} mean? (\textquotedblleft{}A carpenter\textquotedblright{}.) Say it once more.
\end{tcolorbox}
\section[\texorpdfstring{\te{రచయిత} (racayita)}{రచయిత (racayita)}]{\te{రచయిత} (racayita) — a writer, an author}

\label{lesson:TE-C197-racayita}



\begin{quote}
Before the new one: say the Telugu for an application, then the Telugu for a carpenter.
\end{quote}

\subsection*{You'll want to know: \te{రచయిత}}
\textbf{\te{రచయిత}} — \emph{racayita} — \textquotedblleft{}a writer, an author\textquotedblright{}.

\begin{grammarlens}[title={What you've built}]
One more word for this chapter.
\end{grammarlens}

\subsection*{Your turn}
\begin{itemize}
\raggedright
  \item \emph{Say it:} \emph{racayita}
  \item \emph{Say it:} \emph{racayita}, once more
  \item \emph{Say it:} say \emph{racayita}
  \item \emph{Recall:} say the Telugu for an application, then the Telugu for a carpenter, then say \emph{racayita} again
\end{itemize}

\begin{tcolorbox}[breakable,colback=teal!4,colframe=teal!35!black,title={Before you move on}]
What does \te{రచయిత} mean? (\textquotedblleft{}A writer, an author\textquotedblright{}.) Say it once more.
\end{tcolorbox}
\section[\texorpdfstring{\te{సమస్య} (samasya)}{సమస్య (samasya)}]{\te{సమస్య} (samasya) — a problem}

\label{lesson:TE-C197-samasya}



\begin{quote}
Before the new one: say the Telugu for a carpenter, then the Telugu for a writer.
\end{quote}

\subsection*{You'll want to know: \te{సమస్య}}
\textbf{\te{సమస్య}} — \emph{samasya} — \textquotedblleft{}a problem\textquotedblright{}.

\begin{grammarlens}[title={What you've built}]
One more word for this chapter.
\end{grammarlens}

\subsection*{Your turn}
\begin{itemize}
\raggedright
  \item \emph{Say it:} \emph{samasya}
  \item \emph{Say it:} \emph{samasya}, once more
  \item \emph{Say it:} say \emph{samasya}
  \item \emph{Recall:} say the Telugu for a carpenter, then the Telugu for a writer, then say \emph{samasya} again
\end{itemize}

\begin{tcolorbox}[breakable,colback=teal!4,colframe=teal!35!black,title={Before you move on}]
What does \te{సమస్య} mean? (\textquotedblleft{}A problem\textquotedblright{}.) Say it once more.
\end{tcolorbox}
\section[\texorpdfstring{\te{పరిష్కారం} (pariṣkāraṁ)}{పరిష్కారం (pariṣkāraṁ)}]{\te{పరిష్కారం} (pariṣkāraṁ) — a solution}

\label{lesson:TE-C197-pariskaram}



\begin{quote}
Before the new one: say the Telugu for a writer, then the Telugu for a problem.
\end{quote}

\subsection*{You'll want to know: \te{పరిష్కారం}}
\textbf{\te{పరిష్కారం}} — \emph{pariṣkāraṁ} — \textquotedblleft{}a solution\textquotedblright{}.

\begin{grammarlens}[title={What you've built}]
One more word for this chapter.
\end{grammarlens}

\subsection*{Your turn}
\begin{itemize}
\raggedright
  \item \emph{Say it:} \emph{pariṣkāraṁ}
  \item \emph{Say it:} \emph{pariṣkāraṁ}, once more
  \item \emph{Say it:} say \emph{pariṣkāraṁ}
  \item \emph{Recall:} say the Telugu for a writer, then the Telugu for a problem, then say \emph{pariṣkāraṁ} again
\end{itemize}

\begin{tcolorbox}[breakable,colback=teal!4,colframe=teal!35!black,title={Before you move on}]
What does \te{పరిష్కారం} mean? (\textquotedblleft{}A solution\textquotedblright{}.) Say it once more.
\end{tcolorbox}
\section[\texorpdfstring{\te{అభిప్రాయం} (abhiprāyaṁ)}{అభిప్రాయం (abhiprāyaṁ)}]{\te{అభిప్రాయం} (abhiprāyaṁ) — an opinion}

\label{lesson:TE-C197-abhiprayam}



\begin{quote}
Before the new one: say the Telugu for a problem, then the Telugu for a solution.
\end{quote}

\subsection*{You'll want to know: \te{అభిప్రాయం}}
\textbf{\te{అభిప్రాయం}} — \emph{abhiprāyaṁ} — \textquotedblleft{}an opinion\textquotedblright{}.

\begin{grammarlens}[title={What you've built}]
One more word for this chapter.
\end{grammarlens}

\subsection*{Your turn}
\begin{itemize}
\raggedright
  \item \emph{Say it:} \emph{abhiprāyaṁ}
  \item \emph{Say it:} \emph{abhiprāyaṁ}, once more
  \item \emph{Say it:} say \emph{abhiprāyaṁ}
  \item \emph{Recall:} say the Telugu for a problem, then the Telugu for a solution, then say \emph{abhiprāyaṁ} again
\end{itemize}

\begin{tcolorbox}[breakable,colback=teal!4,colframe=teal!35!black,title={Before you move on}]
What does \te{అభిప్రాయం} mean? (\textquotedblleft{}An opinion\textquotedblright{}.) Say it once more.
\end{tcolorbox}
